% m7_data.tex -- Main text, Section 7: confrontation with documented outbreaks, contact records, tracer
% measurements and a school epidemic.
% Paths in "% src:" comments are relative to the repository root.
% Units in this section: hours and metres (D_p, D_air in m^2/h; \remrate in 1/h).
% Scripts (all in code/): s8_outbreaks_checks.py, s8_outbreaks_numbers.py, s8_outbreaks_fig_callcentre.py,
% s8_first_test_exact.py (first test with the round-off-safe propagator), v2_numbers.py and v2_figures.py (second
% set of analyses: every number and generated table body).  Details: Supplementary Sections S6 to S8.
\section{Confrontation with data}\label{s8_outbreaks:sec}\label{s8b_evidence:sec}

Sections~\ref{s3_r0:sec}--\ref{s6_scenes:sec} are statements about a model. This section asks what data say about
it. The data are of four kinds: documented outbreaks (38 records in two datasets; 18 are calibrated on or held out in the two distance-gradient tests, and four
more enter the level checks and negative controls of the first test), RFID
contact records from six indoor settings, tracer measurements in vehicles and rooms, and an influenza epidemic in
29 schools. Five primary tests were run: two out-of-sample tests against outbreaks, and one analysis each of the
contact records, of the tracer measurements and of the zone-level structure. Each follows a protocol fixed
before the model under test was run on its data. We call the protocols \emph{pre-specified} and not
pre-registered, for three reasons that apply to all of them. They are \emph{data-aware}: the outbreak counts,
contact records and school data had been read before the protocols were written, and choices such as strata,
priors and exclusions were made with that knowledge. They are \emph{model-blind}: no output of the model under
test existed when they were frozen (for the tracer analysis, the outcome of the outbreaks under the other models
was already known). And their timing is \emph{self-attested}: they were frozen by SHA-256 hashes of files on the
author's machine, without an external time stamp, version-control history or registry; the protocol files
released with the code are redacted copies (every changed line is listed in the repository), and the hashes refer
to the unredacted originals, so this timing cannot be checked from the repository alone (Supplementary
Section~\ref{S-appC_protocol:sec}). Comparisons added after
the results were known are called \emph{post hoc} and are marked as such; Supplementary
Table~\ref{S-appS_prot2:tab:register} lists every test with its status. Times are in hours and lengths in metres
throughout this section.
% src: data/bsc_outbreaks/outbreaks.csv (27 rows); code/bsc_validation2/a1_more_outbreaks/PREREGISTRATION.md section 2.1 (11 events); events in a test: 2 + 5 (first), 6 + 5 (second); level checks and negative controls of the first test: T3, C3, S1, T7 (outbreaks.csv protocol_role; data/appC_protocol_tables.tex rows A3, A4)
% src: code/bsc_validation/PREREGISTRATION.md section 0; code/bsc_validation/PREREG_FREEZE.json; notes/VERIFICATION.md (Section 4.5, first outbreak test, item 1: partially confirmed); notes/VERIFICATION.md section 4.10 (reservations common to all analyses)

\subsection{Outbreak records}\label{s8_outbreaks:sec:dataset}

The first dataset has 27 records from 31 published sources, each with a registered DOI, 23 of them with counts
of cases and of exposed people: five offices and other workplaces, four retail premises, five classrooms and
class-like rooms, seven vehicle-cabin records, two restaurants, two choir or church events, and two records of
large settings outside the scope of the model (a cruise ship and worker dormitories), recorded for reference as
examples of shared settings that a model of one closed room does not cover; they take part in no comparison
(Supplementary Section~\ref{S-appB_dataset:sec}). We found no documented outbreak in a metro carriage
with a known number of exposed riders (the search was not exhaustive); bus, coach, train and aircraft cabins
stand in for it, and the supermarket has one record with denominators and no geometry.
% src: data/bsc_outbreaks/outbreaks.csv; data/bsc_outbreaks/tables/doi_verification.json; data/appB_dataset_numbers.json -> records_by_group, records_with_counts (23)
The second dataset, collected after the first test, is restricted to events that publish the position of every
exposed person or counts by distance: eight flights
\citep{olsen2003transmission,kenyon1996transmission,baker2010transmission,young2014international,toyokawa2022transmission,speake2020flight,swadi2021genomic,hoehl2020assessment},
a hospital ward observed twice with the same index patient \citep{wong2004cluster,yu2005temporal}, and a
meat-processing line \citep{guenther2020}, which is record O4 of the first dataset with its counts by distance
added. The 38 records of the two datasets therefore describe 35 distinct outbreaks: besides the processing line,
the two records of the ward and those of the coach and the minibus (records T2 and T3) are two exposure cohorts
of one outbreak each.
% src: data/bsc_outbreaks/outbreaks.csv (T3 protocol_role "same index as T2"; O4); data/outbreaks2/SOURCES.md (ward: students and inpatients, same index patient); 38 - 3 = 35
The pathogens are SARS-CoV-2 (five events), SARS-CoV-1 (three),
influenza A(H1N1)pdm09 (two) and \emph{Mycobacterium tuberculosis} (one).
% src: code/bsc_validation2/a1_more_outbreaks/PREREGISTRATION.md section 2; data/outbreaks2/SOURCES.md; data/v2_tab_events2.tex

The records themselves show a distance gradient. With `near' defined as within two rows of a source on aircraft
(the zone of the contact-tracing rule discussed by \citet{hertzberg2016risk}), the index patient's own bay or
cubicle in the ward, and 8\,m on the processing line, the near/far risk ratio pooled over the eleven records of
the second dataset is 3.98 (95\,\% interval 2.32--6.82; random effects, DerSimonian--Laird), with moderate
heterogeneity (Cochran's $Q=18.9$ on 10 degrees of freedom, $p=0.041$; $I^2=0.47$; Supplementary
Fig.~\ref{S-s8_outbreaks:fig:forest}); with fixed effects it is 3.16. The two ward records share the index
patient and are treated here as independent studies; without the second of them the pooled ratio is 4.64
(2.57--8.40). The 8\,m zone is the radius that the source of the processing line singled out as the most
significant on these same data; without the line the pooled ratio is 3.30 (1.99--5.45; post hoc), and over the
eight flights alone 3.75 (2.01--7.00). By pathogen it is 6.66 (3.38--13.1) for SARS-CoV-2, 2.34 (1.53--3.59)
for SARS-CoV-1, and 1.15 (0.35--3.72) on the two influenza flights, which show no detectable gradient. Over the
four held-out events of the first test with distance strata the pooled ratio with the same estimator is 2.05
(1.15--3.65; 2.00 with fixed effects), and these four are mutually compatible (Cochran's $Q=3.77$ on 3 degrees
of freedom, $p=0.29$; $I^2=0.20$). The two pooled values rest on different selections (all eleven records of the
second dataset, calibration events included, against four held-out events of the first) and are not a
comparison of settings. No adjustment was made for companions or households seated together, and seat maps are
published selectively. This is a statement about data, not a test of any model.
% src: data/v2_numbers.json (out2.rr_pooled.*: "all 11 events" RR_random 3.978, RR_fixed 3.159, I2 0.472; "without W2" 4.643 (2.567-8.399); "without P1": 3.296 (1.992-5.455), computed by code/v2_numbers.py; "aircraft (8)" 3.747 (2.006-6.998)); data/bsc_validation2/a1_more_outbreaks/07_descriptive.json (pooled, by pathogen)
% src: data/s8_outbreaks_checks.json (heterogeneity: cochran_Q 3.772, I2 0.205, pooled_rr_random_DL 2.047 (1.147-3.654), pooled_rr 1.996 (fixed effects))

\subsection{Three models and two generic competitors}\label{s8_outbreaks:sec:models}

In a single event of a few hours there is one infectious person and no second generation. All models give the
infection probability of susceptible $j$ in event $k$ as
\begin{equation}\label{s8_outbreaks:eq:risk}
  p_j = 1-\exp(-\varepsilon_k X_j),
\end{equation}
with one infectiousness parameter $\varepsilon_k$ per event and an exposure $X_j$ computed on a rectangular
lattice with reflecting walls, one site per seat position and physical pitches $(a_x,a_y)$; $P_t(y|x)$ is the
propagator of the symmetric walk with hop rates $D/a_x^2$ and $D/a_y^2$. With the index case at $x_0$,
susceptible $j$ at $x_j$ for a time $T_j$, and $\ncell$ sites,
\begin{align}
  \Mzero:&\quad X_j=\frac{1}{\ncell}\left[\frac{T_j}{\remrate}-\frac{1-\e^{-\remrate T_j}}{\remrate^{2}}\right],
     \label{s8_outbreaks:eq:m0}\\
  \Mone:&\quad X_j=\int_0^{T_j}\sum_{x}P_t(x|x_0)\,P_t(x|x_j)\,\dd t=\int_0^{T_j}P_{2t}(x_j|x_0)\,\dd t
     \qquad(D=\Dp),\label{s8_outbreaks:eq:m1}\\
  \Mtwo:&\quad X_j=\int_0^{T_j}(T_j-s)\,\e^{-\remrate s}\,P_s(x_j|x_0)\,\dd s\qquad(D=\Dair).
     \label{s8_outbreaks:eq:m2}
\end{align}
$\Mzero$ is the transient Wells--Riley model of a well-mixed room \citep{wells1955airborne,riley1978,
gammaitoni1997using}. $\Mone$ is the model of Section~\ref{s2_model:sec}: people walk, and transmission occurs
only between two people on the same site (the same-cell kernel), so the exposure is their expected
co-occupation time; the form \eqref{s8_outbreaks:eq:risk} is then exact to first order in $\varepsilon_k$ and,
by Jensen's inequality, an upper bound otherwise. In $\Mtwo$ the index case emits infectious quanta that perform
the lattice walk with coefficient $\Dair$ and are removed at rate $\remrate$, the ventilation rate plus
$0.93\,\mathrm{h}^{-1}$ for deposition and inactivation (the values used by \citet{ou2022} after
\citet{miller2021}); a susceptible inhales in proportion to the quanta on its site. $\Mzero$ is the limit
$\Dair\to\infty$ of $\Mtwo$. $\Mtwo$ belongs to the spatially resolved models of airborne exposure of eddy-diffusion
type \citep{cheng2011modeling,venkatram2021modeling,lau2022predicting}; it shares only the lattice propagator with
$\Mone$ (Table~\ref{s8_outbreaks:tab:models}).
% src: code/bsc_validation/PREREGISTRATION.md section 3; notes/VERIFICATION.md (Section 4.5, first outbreak test, further point 5)

Two generic competitors have no lattice and no physics. The \emph{constant-ratio model} gives every susceptible
in a near zone $\varrho$ times the infection pressure of everyone else; near is the stratum of the source in the
first dataset and, in the second, within two rows of a source on aircraft. In the first test it entered only
after unblinding, with fixed values $\varrho=2$ to 3; in the second it is a pre-specified competitor with one
$\varrho$ per setting class calibrated on the same events as $\Mtwo$. The \emph{exponential kernel} takes
$X_j\propto\exp(-d_j/\ell)$, with $d_j$ the distance to the source and one length $\ell$ per setting class; it was
added post hoc during the re-check of the second test.
% src: code/bsc_validation2/a1_more_outbreaks/PREREGISTRATION.md section 3 (CRR); notes/VERIFICATION.md (Section 4.6, second outbreak test, item 15)

\begin{table}[tbp]
  \centering\small
  \caption{The three models compared. Each has one infectiousness parameter $\varepsilon_k$ per event. The
  mixing coefficient is shared by all events of a setting class (first test: vehicle cabins, seated rooms;
  second test: aircraft cabins, rooms); the removal rate $\remrate$ is measured or integrated over a prior
  and is never fitted.}
  \label{s8_outbreaks:tab:models}
  \begin{tabular}{@{}l>{\raggedright\arraybackslash}p{0.36\textwidth}>{\raggedright\arraybackslash}p{0.2\textwidth}>{\raggedright\arraybackslash}p{0.22\textwidth}@{}}
    \toprule
    Model & Transmission & What walks on the lattice & Shared parameter \\
    \midrule
    $\Mzero$ & well-mixed air, eq.~\eqref{s8_outbreaks:eq:m0} & nothing & none \\
    $\Mone$  & between two people on the same site, eq.~\eqref{s8_outbreaks:eq:m1} & people & $\Dp$ \\
    $\Mtwo$  & inhalation of quanta emitted by the index case, eq.~\eqref{s8_outbreaks:eq:m2} & quanta & $\Dair$ \\
    \bottomrule
  \end{tabular}
\end{table}

\subsection{Protocols of the two outbreak tests}\label{s8_outbreaks:sec:protocol}

\emph{First test.} Vehicle cabins are calibrated on a cohort of contacts of 2{,}334 index cases on high-speed
trains \citep{hu2021} (attack rates in 22 seat-offset cells, the adjacent-seat cell excluded because travelling
companions inflate it) and seated rooms on the Guangzhou restaurant \citep{lu2020,li2021} (interval-censored
counts, because only one infection per neighbouring family is certain to have occurred at the lunch). Four
single events (a bus \citep{shen2020}, a coach \citep{luo2020,ou2022}, a flight \citep{khanh2020} and a
classroom \citep{lamhine2021}) are held out and tested on the split of their cases
between a near and a far stratum given the event total, and the Seoul call centre \citep{park2020} on the
join-count $z$ of its case seats. The score $S$ is the logarithm of the predictive probability of the observed
split and $\Delta=S(\text{model})-S(\Mzero)$. The decision rule requires the pooled $\Delta$ to be positive with
one-sided $p<0.05$ under $\Mzero$ (A1), each event to have a two-sided predictive $p\ge0.05$ (A2), the case
counts of two events transferred from the calibration settings (a minibus \citep{luo2020} and masked
classrooms \citep{hershow2021}) to lie inside
their 90\,\% predictive intervals (level checks, A3), and no over-prediction in four groups with few or no cases
(negative controls, A4); two or more failures of A2 give outcome W4, a list of what is and is not
reproduced and no general statement of agreement. The protocol was frozen 87 seconds before the first model run
(Supplementary Section~\ref{S-appC_protocol:sec}).
% src: code/bsc_validation/PREREGISTRATION.md sections 2.1 (V1, V2), 3.4, 5.1, 5.2, 6, 8; code/bsc_validation/PREREG_FREEZE.json

\emph{Second test.} It was designed around the weakness of the first, which a constant ratio passes
(Section~\ref{s8_outbreaks:sec:holdout}). The split does not use outcomes: the six events caused by pathogens
other than SARS-CoV-2, all published before 2020, calibrate; the five SARS-CoV-2 events (four flights and the
processing line) are held out. The observable is the vector of case counts over two to four distance strata,
given the event total; $\Delta_0$ and $\Delta_{\mathrm{C}}$ are the summed log-score differences of $\Mtwo$
against $\Mzero$ and against the constant-ratio model. The rule: B1, $\Delta_0>0$ with one-sided $p<0.05$ under
$\Mzero$; B2, $\Delta_{\mathrm{C}}>0$ with one-sided $p<0.05$ under the constant-ratio model (the mirror
statement, $\Delta_{\mathrm{C}}<0$ with $p<0.05$ under $\Mtwo$, means that the constant ratio predicts better);
B3, every held-out event has a two-sided predictive $p\ge0.05$; two or more failures of B3 give the outcome V5, a
list of the events reproduced and of those not. Computed after calibration and before any held-out stratum count
was scored, the power was (with the frozen propagator; the corrected one described below gives the values in
brackets): if the calibrated $\Mtwo$ were true, B1 and B2 would both be met with probability 0.96 (0.93; 0.82
on the four flights alone); if outbreaks followed a constant ratio, with probability 0.03 (0.05), while B1 alone
would be passed with probability 0.65 (0.73), so that only B2 can count as evidence for the kernel.
% src: code/bsc_validation2/a1_more_outbreaks/PREREGISTRATION.md sections 0, 4-7; code/bsc_validation2/a1_more_outbreaks/PREREG_ADDENDUM_power.md; data/v2_numbers.json (out2.power.registered: M2 P_B1_B2 0.956, CRR P_B1_B2 0.028, CRR P_B1 0.652; out2.power.registered_air: M2 P_B1_B2 0.821; out2.power.corrected: M2 B1B2 0.926, CRR B1B2 0.049, CRR B1 0.732)

\emph{One deviation from the frozen code} affects both tests. The frozen propagator is a sum over cosine modes
whose absolute error reaches about $10^{-10}$ of the largest exposure in an event; at small mixing coefficients
the true exposures of distant receivers are smaller than that, and the mode sum returns round-off noise of
either sign in their place, which makes far infections look possible. This was found after the held-out events
of the second test had been scored. All results below use a propagator that keeps the mode sum where it is
resolved and replaces it elsewhere by a positive sum of image terms; it agrees with a third method
(uniformisation) to better than $10^{-8}$ relative in the classroom of the first test, and with the image sum to
about $10^{-4}$ in the ward of the second, where both are resolved. The correction changes no number for the flights and no
outcome of either rule; the values obtained with the frozen propagator are reported in Supplementary
Sections~\ref{S-s8_outbreaks:sec:holdoutdetail} and~\ref{S-appS_prot2:sec:numerics}.
% src: data/s8_first_test_exact.json (kernel.rows: spectral_min_rel down to -1.28e-10, fixed_vs_uniformisation_max_rel_err <= 8.4e-9, classroom C1 layouts); code/bsc_validation2/a1_more_outbreaks/POSTHOC_LOG.md (P5); data/bsc_validation2/a1_more_outbreaks/09_numerics_check.json (largest relative difference 8.4e-5, ward inpatients; Supplementary Section S10)

\subsection{First test}\label{s8_outbreaks:sec:holdout}

On the train cohort both spatial models gain 26.6--26.7 units of log-likelihood over $\Mzero$ for one more
parameter; $\Mtwo$ gives $\Dair=25\sqmh$ (90\,\% posterior interval 9.7--43; here and below the posterior is a
pseudo-posterior proportional to the profile likelihood under a flat prior on $\log_{10}D$, Supplementary
Section~\ref{S-s8_outbreaks:sec:calibration}) and $\Mone$ gives $\Dp=0.50\sqmh$
(0.31--0.79). Both are nevertheless rejected by the data they were calibrated on (deviance 72.6--72.9 on 20
degrees of freedom): the isotropic kernel cannot follow the anisotropy of the train table. For the restaurant
$D$ is identified only within three decades, and the fitted $\Mtwo$ conflicts with the tracer-gas measurement of
\citet{li2021} (Supplementary Section~\ref{S-s8_outbreaks:sec:calibration}).
% src: data/bsc_validation/tables.md (Table V1); data/s8_outbreaks_numbers.json (calibration.trains.M2: joint_max_D 25.1, q05 9.67, q95 43.4)

\begin{table}[tbp]
  \centering\small
  \caption{First test, hold-out events: near/far split (event totals fixed, one $\varepsilon_k$ per event). RR:
  risk ratio near/far, observed with its 95\,\% interval, or the ratio of the expected stratum attack rates of a
  model; $p$: two-sided predictive $p$-value; $\Delta$: log-score difference to $\Mzero$. $\Mone$ is evaluated
  in the exact simulator for the bus, the flight and the classroom. Sources: bus \citep{shen2020}, coach
  \citep{luo2020,ou2022}, flight \citep{khanh2020}, classroom \citep{lamhine2021}.}
  \label{s8_outbreaks:tab:holdout}
  % src: data/s8_first_test_exact.json (primary.events.*.{M0,M1,M2}: rr_pred, p_two_sided, delta_vs_M0; primary.pooled.M2.all: delta 2.835, p_under_M0 2.48e-4; without_T5: -3.195; primary.pooled.M1.all: -2.939); data/bsc_validation/tables.md (Tables V5, V6)
  \setlength{\tabcolsep}{3.4pt}
  \begin{tabular}{@{}lccccccc@{}}
    \toprule
    & & & $\Mzero$ & \multicolumn{2}{c}{$\Mone$} & \multicolumn{2}{c}{$\Mtwo$} \\
    \cmidrule(lr){4-4}\cmidrule(lr){5-6}\cmidrule(l){7-8}
    Event (near stratum) & Near vs far & Observed RR & $p$ & RR; $p$ & $\Delta$ & RR; $p$ & $\Delta$ \\
    \midrule
    Bus (rows 5--11)         & 14/33 vs 9/34 & 1.60 (0.81--3.19) & 0.20   & 4.0; n.a.$^{a}$ & $+0.30$ & 2.4; 0.38  & $+0.32$ \\
    Coach (rear rows)        & 3/19 vs 4/26  & 1.03 (0.26--4.06) & 1.00   & 34; 0.0002      & $-7.64$ & 11; 0.007  & $-3.94$ \\
    Flight ($\le2$ seats away) & 11/12 vs 1/8  & 7.33 (1.16--46.2) & 0.0008 & 1.4; 0.009      & $+2.38$ & 3.6; 0.61  & $+6.03$ \\
    Classroom (front rows)   & 8/10 vs 4/14  & 2.80 (1.16--6.78) & 0.036  & 1.4; 0.31$^{b}$ & $+2.02$ & 5.5; 0.068$^{c}$ & $+0.43$ \\
    \midrule
    Pooled & & & & & $-2.94$ & & $+2.83^{d}$ \\
    \bottomrule
  \end{tabular}

  \smallskip
  \begin{minipage}{0.96\textwidth}\footnotesize
  $^{a}$\,Observed total unattainable: $\Prob(\text{total}\ge23)=0.0002$; conditional on the total, $p=0.46$.
  $^{b}$\,Only 2 of 24 posterior draws can reach 12 cases.
  $^{c}$\,Not a genuine prediction: rejected at the posterior mode ($p=0.0009$); see text.
  $^{d}$\,Exact one-sided $p=2.5\times10^{-4}$ under $\Mzero$; $-3.19$ without the flight.
  \end{minipage}
\end{table}

Table~\ref{s8_outbreaks:tab:holdout} gives the hold-out results. For $\Mtwo$: \emph{Flight VN54} is reproduced
(11 of the 12 passengers seated at most two seats from the index case and one of the 8 farther away were
infected; $\Mzero$ is rejected, $p=0.0008$, and $\Mtwo$ is not, $p=0.61$), robustly to the calibration and to the
ventilation prior; this is the one robust positive result of the first test, and it is thin: 20 exposed
passengers, among them four companion couples seated side by side. The \emph{bus} is compatible but not
discriminating. The \emph{Hunan coach} is a failure: infection was uniform along the coach (3/19 rear, 4/26
front), and $\Mtwo$ puts six of the seven cases in the rear half (predicted ratio 11, $p=0.007$), in 10 of the 12
pre-specified sensitivity variants. The \emph{Marin classroom} is a pass that is not a prediction: the
predictive $p$ of 0.068 comes from the diffuse upper tail of the restaurant posterior, where the kernel is
nearly flat across the room; at the posterior mode the observation is rejected ($p=0.0009$), and in 83\,\% of
the posterior draws the emission rate needed to produce the 12 cases exceeds $10^{4}\qph$. With the frozen
propagator this $p$ was 0.11 (Section~\ref{s8_outbreaks:sec:protocol}).
% src: data/bsc_validation/03_holdout_shape.json (T5, T1, T2: p = 0.0069, k_pred_mean 6.2); data/s8_outbreaks_numbers.json (sensitivity_M2: n_coach_fails 10)
% src: data/s8_first_test_exact.json (mode: D 0.501, p_two_sided 0.00093; vent.primary.events.C1: emission_share_above_1e4 0.833; primary.events.C1.M2.p_two_sided 0.068); data/bsc_validation/03_holdout_shape.json (C1: 0.115)
The \emph{Seoul call centre} is a failure as well. On its published seat map the 79 case seats among the 137
desks of the affected wing show no clustering at desk scale (join-count $z=-0.64$; random relabelling, one-sided
$p=0.76$), whereas the restaurant-calibrated $\Mtwo$ predicts strong clustering (median $z\approx11$); integrated
over the posterior, the probability of a $z$ as small as observed is 0.007 or 0.013, depending on how posterior
draws are weighted, below the limit 0.025 of the central 95\,\% interval. $\Mzero$ is compatible ($P=0.26$).
The protocol estimated this probability from 40 posterior draws (0.031 and 0.017, on either side of the limit);
the quadrature that replaces the estimate was added in the re-check after the results were known, and its value
depends on the upper limit of the flat prior on $\log_{10}D$: 0.007 and 0.013 with the limit $10^{4}\sqmh$ used
throughout, 0.013 and 0.024 if the prior is extended to $10^{6}\sqmh$ (post hoc).
% src: data/bsc_outbreaks/tables/analysis_results.json (callcentre.joincount); data/validation_checks/callcentre_quadrature.json (equal_weight 0.0068, pooled_weight 0.0128; prior_upper_limit_sensitivity 1e6_flat_extension: equal 0.0127, pooled 0.0238); data/bsc_validation/tables.md (Table V9: M2 median z 10.99, M0 0.262; 40-draw estimate 0.031 pooled, 0.017 equal weight)

The pooled score of $\Mtwo$ is $\Delta=+2.83$ (exact one-sided $p=2.5\times10^{-4}$ under $\Mzero$), so A1 is
passed; it rests on the flight, without which $\Delta=-3.19$. Two adequacy failures give the outcome W4 for
$\Mtwo$. And the rule can be passed without any lattice: a model with no calibration that assigns the same
near/far risk ratio of 2, 2.5 or 3 to every event scores $\Delta=+7.6$, $+7.7$ and $+7.1$, higher than $\Mtwo$,
with no single-event failure (a post-hoc comparison). Passing A1 shows that a proximity gradient exists in these
data; it does not show that the lattice kernel or its calibration adds predictive value.
% src: data/s8_first_test_exact.json (primary.pooled.M2: all 2.835, p 2.48e-4; without_T5 -3.195); data/validation_checks/v05_generic.txt (RR = 2: +7.63; 2.5: +7.74; 3: +7.09; failures [])

\emph{The same-site model $\Mone$} fails the pooled test ($\Delta=-2.94$) and three single events, and its
outcome is W4 as well; the failure on the bus rests on rule P4 of the post-hoc log, adopted after the results,
which counts an event whose observed total the exact simulator cannot reach as a failure (Supplementary
Section~\ref{S-appC_protocol:sec}). Calibrated on the trains, $\Dp=0.50\sqmh$: a seated passenger changes site a few times
per hour. At such mobilities the mean number of cases stays below the observed number in three of the four
held-out events however large the transmission rate, because the attack rate saturates at the probability that
a susceptible ever shares a site with the index case: on the bus the mean is at most about 9 cases (23
observed), and the observed totals are reached with probability 0.0002 on the bus, 0.025 on the flight and 0.081
in the classroom (Supplementary Table~\ref{S-s8_outbreaks:tab:m1}). A larger mobility would lift this ceiling but
is excluded by the train data, and at walking mobility the room is well mixed and $\Mone$ loses its spatial
structure.
% src: data/bsc_validation/02b_m1_exact.json (events.*.rows: p_total_ge_K, mean 0.00019 (T1), 0.0247 (T5), 0.0810 (C1)); data/bsc_validation/tables.md (Table V8: 8.9 (7.5-12.2)); data/bsc_validation/03_holdout_shape.json (pooled.M1)

\subsection{Second test}\label{s8_outbreaks:sec:second}

\begin{table}[tbp]
  \centering\footnotesize
  \setlength{\tabcolsep}{3.2pt}
  \caption{Second test: held-out events (event totals fixed, one $\varepsilon_k$ per event; kernel and ratio
  calibrated on the six pre-2020 events). Expected numbers of cases per stratum under the constant-ratio model
  (CR) and under $\Mtwo$; two-sided predictive $p$ of the observed vector under $\Mzero$, CR and $\Mtwo$;
  $\Delta_0$ and $\Delta_{\mathrm{C}}$: log score of $\Mtwo$ minus that of $\Mzero$ and of CR.}
  \label{s8_outbreaks:tab:second}
  % generated: data/v2_tab_holdout2.tex (code/v2_numbers.py <- data/bsc_validation2/a1_more_outbreaks/10_corrected.json, primary_corrected)
  \begin{tabular}{@{}>{\raggedright\arraybackslash}p{0.155\textwidth}c>{\raggedright\arraybackslash}p{0.135\textwidth}>{\raggedright\arraybackslash}p{0.1\textwidth}>{\raggedright\arraybackslash}p{0.1\textwidth}>{\raggedright\arraybackslash}p{0.15\textwidth}cc@{}}
    \toprule
    Event & Near vs far & Observed by stratum (at risk) & Expected, CR & Expected, $\Mtwo$ &
      $p$: $\Mzero$ / CR / $\Mtwo$ & $\Delta_0$ & $\Delta_{\mathrm{C}}$ \\
    \midrule
    \tabinput{data/v2_tab_holdout2}
    \bottomrule
  \end{tabular}
\end{table}

\emph{Calibration.} On the four pre-2020 flights the constant ratio is $\varrho=3.2$ (90\,\% posterior interval
1.6--5.3) and the mixing coefficient of $\Mtwo$ is $\Dair=126\sqmh$ (52--438). In the ward the ratio is 3.5
(1.8--5.5) and the coefficient $794\sqmh$ (302--4{,}390), close to the well-mixed limit, because the two ward
datasets contradict each other: the students, exposed for 40 minutes, ask for a steep kernel and the inpatients,
exposed for days, for a flat one.
% src: data/v2_numbers.json (out2.cal.air|CRR: map 3.16, q05 1.61, q95 5.32; out2.cal.air|M2: 125.9, 51.6, 437.5; out2.cal.room|CRR: 3.55, 1.83, 5.53; out2.cal.room|M2: 794.3, 301.9, 4390.8)

\emph{Against the well-mixed model: passed.} Over the five held-out events $\Delta_0=+14.35$ (no value as large
in 200{,}000 draws under $\Mzero$; $+14.30$ in the re-check, with its own geometry code and propagator), and
$+10.19$ on the four flights alone; every held-out event contributes (Table~\ref{s8_outbreaks:tab:second}).
% src: data/v2_numbers.json (out2.pooled.M2.all: delta0 14.351; out2.pooled.M2.air: 10.193; out2.pooled_check.M2.all: 14.299)

\emph{Against the constant ratio: not passed.} Pooled, $\Delta_{\mathrm{C}}=-9.15$, a value that $\Mtwo$ did not
produce in 200{,}000 draws if it were true: by the mirror criterion the constant ratio predicts better.
Nine-tenths of it ($-8.38$) is the processing line, where the comparison is not between equals, because the near
zone of 8\,m is the radius that the source singled out on these same data; with a zone of 4\,m $\Mtwo$ is ahead on
the line by 2.2, and with 12\,m it trails by 0.6 (post hoc). On the four held-out flights the two models are tied:
$\Delta_{\mathrm{C}}=-0.77$, with $p=0.16$ for `$\Mtwo$ better' and $p=0.065$ for `constant ratio better', in a
test that had a power of 0.82 to show $\Mtwo$ better if it were true. Over the 15 pre-specified sensitivity
variants $\Delta_{\mathrm{C}}$ on the flights runs from $-3.42$ to $+1.21$: the sign depends on definitions of
outcome, strata and near zone that the data cannot fix.
% src: data/v2_numbers.json (out2.pooled.M2.all: deltaC -9.150, p_mirror 0; out2.pooled.M2.room: deltaC -8.378; out2.pooled.M2.air: deltaC -0.772, p_B2 0.160, p_mirror 0.065; out2.pooled.plant_share_of_deltaC 0.916; out2.check.near1_4m: room deltaC +2.19; out2.check.near5_12m: room -0.64; out2.variants.summary: dC_air_min -3.42, dC_air_max 1.21)

\emph{Adequacy: two failures.} Two of the five held-out events lie outside the predictive distribution of
$\Mtwo$: the flight to Naha ($p=0.047$, marginal and fragile; $p=0.25$ if households are counted once) and the
processing line ($p=1.4\times10^{-5}$), where the attack rate is flat out to 8\,m and then drops, whereas the
ward-calibrated kernel is almost well mixed. The outcome is V5. The competitors fail as well: the constant ratio
on one flight ($p=0.020$) and the well-mixed model on three events.
% src: data/v2_numbers.json (out2.hold.F5: M2 p 0.0472; out2.hold.P1: M2 p 1.42e-5; out2.hold.F7: CRR p 0.0195; out2.pooled.M2.decision: B3_failures [F5, P1], outcome V5); data/checks/outbreaks2.json: naha (p 0.25 with households collapsed)

\emph{The same-site model} was scored in the second test in a closed form only, and its protocol fixed the
reading in advance: the first test had shown with the exact simulator that the closed form overstates the reach
of $\Mone$ (pooled $-11.11$ against $-2.94$ for the simulated model), so a pass of the closed form would not bear
on $\Mone$ and a failure would confirm the first test. The closed form passes against the well-mixed model
($\Delta_0=+15.45$, largely through the processing line), is worse than the constant ratio pooled
($\Delta_{\mathrm{C}}=-8.06$) and on the flights ($-4.63$), and fails on three of the five held-out events.
% src: data/v2_numbers.json (out2.pooled.M1.all: delta0 15.446, deltaC -8.055; out2.pooled.M1.air: deltaC -4.634; out2.pooled.M1.decision: B3_failures [F5, F7, P1], V5); code/bsc_validation2/a1_more_outbreaks/PREREGISTRATION.md section 3 (M1cf); data/s8_first_test_exact.json (primary.pooled: M1cf -11.109, M1 -2.939)

\emph{A physics-free kernel does as well.} In a post-hoc comparison the exponential kernel, calibrated in the
same way (decay length on aircraft 2.5\,m), scores $\Delta_0=+18.0$; $\Mtwo$ is ahead of it by 0.31 on the
flights, which is not significant, and behind it by 3.7 over the five events. Nothing in this test distinguishes
the lattice propagator, or the removal of aerosol by ventilation, from a generic smooth decay with distance.
Further pre-specified secondary analyses, among them the transfer of the train kernel of the first test to the
eight flights (it beats the well-mixed model and is rejected on four flights), are in Supplementary
Section~\ref{S-s8_outbreaks:sec:seconddetail}.
% src: data/v2_numbers.json (out2.check.exp_vs_crr.all: delta0 18.02; out2.check.m2_vs_exp.air: deltaC +0.313; out2.check.m2_vs_exp.all: deltaC -3.72; out2.sec.train_kernel.flights: delta0 7.555; out2.sec.train_kernel.padeq); data/checks/outbreaks2.json: exponential_kernel.aircraft_decay_length_m (2.5)

\subsection{Mixing coefficients and levels}\label{s8_outbreaks:sec:mixing}\label{s8_outbreaks:sec:levels}

Fitted alone, the vehicle events of the first test ask for mixing coefficients that differ by orders of
magnitude: 0.06--$0.1\sqmh$ for the flight, 25 for the trains, 63 for the bus and at least $6{,}300\sqmh$ for the
coach; a common value is rejected (likelihood ratio 9.52 on 3 degrees of freedom, approximate $p=0.023$, a post-hoc
test). In the second test a common coefficient is rejected for the three room events (likelihood ratio 17.4 on
2 degrees of freedom, approximate $p=1.6\times10^{-4}$) and, more weakly, for the eight flights (15.2 on 7,
approximate $p=0.034$); these approximate $p$-values use a chi-square reference whose conditions are not met when
maxima lie on the edge of the grid. No setting class has a single coefficient, and none transfers between
classes. One infectiousness parameter is fitted per event, in both tests, so no attack rate in this section is a
prediction; the emission rates implied by five single-index events and one cluster with several instructors
of the first dataset span a factor of 76 (53 if only the confirmed choir cases are counted;
Supplementary Table~\ref{S-s8_outbreaks:tab:levels}). Published outbreaks are selected for being large, and
their level cannot be predicted from a typical index case.
% src: data/appC_protocol_numbers.json (loo.vehicle: LR 9.52, p 0.0231); data/bsc_validation/07_loo.json (classes.vehicle.M2: D_own_mode 25.1, 63.1, 6310, 0.1); data/validation_checks/v10_c1_loo.txt (T5 0.063)
% src: data/v2_numbers.json (out2.D.room: LR 17.44, df 2, p 1.63e-4; out2.D.air: LR 15.18, df 7, p 0.0338)
% src: data/bsc_validation/04_level_controls.json (implied.H1.E_implied = 2734.3; implied.T2.E_implied = 36.0; implied.R2.E_implied = 1921.0); ratios 2734.3/36.0 = 76, 1921.0/36.0 = 53

\subsection{Recorded contacts}\label{s8b_evidence:sec:contacts}

Does the movement mechanism of the model, a random walk with contact on a shared site, describe who meets whom
indoors? Six open SocioPatterns datasets of face-to-face proximity recorded with wearable RFID sensors at a
resolution of 20 seconds \citep{cattuto2010dynamics} were used: an office building in 2013 and in 2015
\citep{genois2015data,genois2018can}, a primary school \citep{stehle2011high}, a hospital ward
\citep{vanhems2013estimating}, a high school \citep{mastrandrea2015contact} and a scientific conference
\citep{genois2018can,stehle2011simulation}, with 75 to 403 people over 2 to 10 days. Each model is calibrated on
the first half of the days of a dataset and predicts the second half. The walk of Section~\ref{s2_model:sec}, in
discrete time with one site per contact range, has two parameters (the number of sites and the hop
probability), fixed by the contact time per jointly present pair and by the mean contact duration. Its
competitors are a well-mixed model and a model with a constant ratio of contact within and between the recorded
groups. Seven families of contact statistics are scored, and an SIR epidemic is run on the recorded and on the
model-generated contacts at the same total contact time in four scenarios; the movement component is validated on
a dataset if it meets the epidemic criterion and at least six of the seven families
(Supplementary Section~\ref{S-appS_prot2:sec:contacts}).
% src: code/bsc_validation2/a2_contact_mobility/PREREGISTRATION.md sections 2-7; data/v2_numbers.json (con.dataset.*: N 75-403, n_days 2-10)

The walk fails on all six datasets (Supplementary Table~\ref{S-s8b_evidence:tab:contacts}). It passes one or two of the seven
families on each, namely those it was calibrated on, and does not reproduce who meets whom: people in the walk
have too many partners on five of the six datasets (6.0 per person and day against 4.6 in the office of 2013, 90
against 47 in the primary school), and contact within the recorded groups is 10 to 33\pct{} of contact time in
the walk and 57 to 94\pct{} in the data. Because the walk spreads the same contact time over too many, too equal
partners, the index case infects more people on the walk's contacts than on the recorded ones in all 24
scenarios, by 25 to 72\pct{} on average per setting. The two models without space do neither better nor worse,
and the outcome is the lowest of the protocol's ladder of outcomes (A0 in its labels), with the sentence that the
protocol reserved for this case: the lattice walk adds nothing over a non-spatial contact model. A static network in which each pair keeps the contact rate it had on the
training days does better (post hoc). In the walk the number of contacts grows as the square of the number of
people present; in four of the five settings where this can be estimated it grows at most linearly (exponents
0.45 to 0.95), and in the hospital ward faster than the square (2.8).
% src: data/v2_tab_datasets.tex (S7_alpha observed / walk: 0.45/1.97, 2.83/2.20, 0.95/2.06, 0.76/2.12, 0.81/1.95)
% src: data/v2_numbers.json (con.model.RW0.S: [1,2,2,1,1,1]; con.stats_obs_vs_RW0: S3_deg_mean 4.58/5.96 (InVS13), 46.6/90.4 (LyonSchool); S6_within_frac observed 0.57-0.94, walk 0.10-0.33; con.model.RW0: n_cells_over 24 of 24, ratio_min 1.253, ratio_max 1.715; con.model.WM; con.ladder: A0; con.check.static); data/v2_tab_datasets.tex (S7_alpha)
Walks anchored to home sites, developed afterwards on three datasets and frozen before three confirmatory ones,
reduce the error of simulated attack rates from 0.17 to 0.04--0.05 in a description after the fact, but their
confirmatory test is uninformative, because the frozen fitting procedure does not recover its own synthetic
truth; anchoring is a direction, not a result. A walk with two zones of different mobility, the only variant
that touches heterogeneous mobility, passes two of the seven families on each dataset and does not improve on the
homogeneous walk (Supplementary Section~\ref{S-s8b_evidence:sec:contactsdetail}). The calibrated hop probabilities correspond to about 220 to 450
sites$^2$ per day, between $\Dz=1$, which was not tested against contacts, and walking mobility.
% src: data/v2_numbers.json (con.model.RWhet.S: [2,2,2,2,2,2]; con.model.RWhet.E_datasets 0; con.model.RW0.mean_abs_dA 0.167; con.model.RWclus: 0.041; con.model.RWstickZ: 0.038; con.model.RWstickO: 0.050; con.check.self_power); per-direction hop rate (p/4) x 4,320 steps per day = 219-451 (con.dataset.*.p: 0.203-0.417)

\subsection{The airborne kernel fixed by tracer measurements}\label{s8b_evidence:sec:tracer}

If the non-local term of $\Mtwo$ describes the transport of aerosol, its mixing coefficient can be measured
without any outbreak. Published tracer measurements exist for four of the settings: ethane at eight seats of the
Hunan coach \citep{ou2022}, particles released from a breathing mannequin in a wide-body aircraft cabin
\citep{kinahan2021aerosol}, salt aerosol in a rail carriage \citep{woodward2022evaluation}, and tracer gas at
seven tables of the Guangzhou restaurant \citep{li2021}; for the classroom a regression of the turbulent diffusion
coefficient on the air change rate measured in two houses \citep{cheng2011modeling} is extrapolated. In the three
vehicles the measured fields, summarised by a diffusion coefficient, are of order 100 to $1{,}000\sqmh$: about
$1{,}300$ in the coach (bounded from below only), 700 and 110 for two releases in the cabin, and 900 and 350 in
the carriage. These coefficients are points of a grid spaced by a factor of 1.12 (those fitted to outbreaks in
both tests, of a grid spaced by 1.26), so that only their first two digits carry information. They are 4.5 to about 50 times the
vehicle value of the first test ($25\sqmh$), and the extrapolated room values (169 to $185\sqmh$) a few hundred
times its room value ($0.5\sqmh$). The aircraft coefficient of the second test
($126\sqmh$) lies inside the range measured in cabins; that does not make an outbreak-fitted coefficient a
measurement of air mixing, since within the same test no common coefficient holds. The coefficient that
outbreaks ask for absorbs whatever makes risk fall steeply with distance: exposure at close range, companions
seated together, behaviour.
% src: code/bsc_validation2/a3_tracer_physics/PREREGISTRATION.md section 3; data/v2_numbers.json (tr.D.coach: D 1259, lo 708; tr.D.cabin: D_5L 708, D_5A 112; tr.D.train: middle 891, end 355; tr.D.table: classroom 169, office 185; out2.cal.air|M2: 126); data/checks/tracer.json: ratio_tracer_to_outbreak_D_at_low_cabin_value (4.5); grids: code/bsc_validation2/a3_tracer_physics/PREREGISTRATION.md section 3 (log10 step 0.05, factor 1.122; 1258.93 = 10^3.10, 707.95 = 10^2.85, 112.20 = 10^2.05, 891.25 = 10^2.95, 354.81 = 10^2.55), code/bsc_validation/PREREGISTRATION.md and code/bsc_validation2/a1_more_outbreaks/PREREGISTRATION.md (61 points over [0.01, 10^4], factor 1.259); tr.D.table: classroom 169.16, office 185.2

With the kernel fixed in this way, the events of the first dataset were predicted with one free parameter each
(Supplementary Table~\ref{S-s8b_evidence:tab:tracer}). The coach and the call centre are no longer rejected, but only because the
measured kernel is close to well mixed, and the well-mixed model fits as well or better. The classroom is the one
gain over the well-mixed model ($p=0.42$ against $0.036$), matched by constant ratios of 2 and 3. Flight VN54 is
not reproduced: the predicted ratio is 1.66 against 7.3 observed, and the formal non-rejection ($p=0.14$) depends
on zero sensor readings (with the zeros treated as missing, a variant listed in the protocol, $p=0.003$). The
restaurant is not reproduced at the reported count, and the train table is a failure for every model. Over the
five stratum events the tracer kernel scores $+7.87$ against the well-mixed model and $-0.80$ and $-0.78$ against
the constant ratios 2 and 3; the outcome of the tracer protocol is the one for one adequacy failure and no gain over a generic gradient
(S4 in its labels), and the one for two adequacy failures (S5) if the flight, whose formal pass rests on zero
sensor readings, is counted as a failure. The tracer kernel is not
validated as a spatial predictor. In an exploratory extension to the four held-out flights of the second test,
written after its results were known, the two cabin coefficients beat the well-mixed model and tie with a
constant ratio of 3 ($-0.16$; Supplementary Section~\ref{S-s8b_evidence:sec:tracerdetail}).
% src: data/v2_numbers.json (tr.event.C1: P p 0.424, M0 p 0.036; tr.event.T5: P p 0.136, rr_expected 1.660, KA_zeros_missing 0.0028; tr.eval.pooled.P: M0 delta 7.873; CRR2 -0.800; CRR3 -0.780; tr.eval.decision: outcome S4); notes/VERIFICATION.md (Section 4.8, tracer measurements, item 11: outcome S5 if the flight is counted as a failure); data/v2_numbers.json (tr.ext.pooled.P: M0 +10.36, CRR3 -0.16)

\subsection{Zone-level structure}\label{s8b_evidence:sec:zones}

At the scale of zones the model says that transmission between two groups is proportional to the time their
members spend in the same place, with the frequency-dependent normalisation of Section~\ref{s2_model:sec}: if a
person of group $i$ spends a fraction $f_{iz}$ of the time in zone $z$, the force of infection on group $i$ is
\begin{equation}\label{s8b_evidence:eq:zone}
  \lambda_i=\beta\sum_z q_z\,f_{iz}\,\frac{\sum_jf_{jz}I_j}{\sum_jf_{jz}N_j},
\end{equation}
a residence-time formulation in the sense of \citet{bichara2015sis}. The model fitted puts measured shares of
contact duration on three exclusive levels (the own class, the other classes of the grade, the rest of the
school), which identifies contact time measured by RFID with time spent together. The epidemic is seasonal
influenza in the 29 public primary schools of Matsumoto, Japan, in 2014/15, with the day of onset of 2{,}548
infected among 10{,}923 pupils in 455 classes \citep{endo2021within}; the mixing is the share of a pupil's contact
time measured in a primary school in Lyon \citep{stehle2011high,gemmetto2014mitigation}: 0.762 with the own
class, 0.104 with the other classes of the grade and 0.134 with the rest of the school. A daily chain-binomial
model with the serial-interval kernel of the source is fitted under two-fold cross-validation by school; the rule
has three tiers (Supplementary Section~\ref{S-appS_prot2:sec:zones}).
% src: code/bsc_validation2/a4_zone_level/PREREGISTRATION.md sections 1, 2.1-2.6; data/v2_numbers.json (zone.counts: 10923, 2548, 455, 29); data/schools/lyon_mixing.json (shares 0.762, 0.104, 0.134)

The zone model predicts held-out schools far better than a well-mixed school (ahead in 27 of 29 schools) and
than independent classes (Supplementary Table~\ref{S-s8b_evidence:tab:zone}). The share between classes of the same grade agrees
to two decimals; exact agreement of the three shares is rejected, the fitted school-level share being almost
twice the measured one. This comparison is confounded: the model has no household term, whereas the source
attributes 41\pct{} of the pupils' infections to households, and siblings attend the same school in different
grades. The fitted shares lie within the tolerance, but not securely (11\pct{} of bootstrap resamples exceed
it). The pre-specified generic two-level model with one fitted ratio is not beaten, and a guess of 0.7, 0.1 and
0.2, chosen after the fitted shares were known, does better. The verdict of the rule, `right structure within
tolerance, not exact', is reported as a weak pass: the measurement is not shown to matter beyond the statement
that most transmission is within the class. The test rests on the residence-time reduction
\eqref{s8b_evidence:eq:zone}, which is an assumption; the lattice walk, the gradient within a room, the airborne
term and attack-rate levels are not involved. A zone model of the Seoul call-centre building by floor excludes a
well-mixed building, which needs no model, and nothing more: the time spent in the shared lifts and lobby is not
measured, and with a longer assumed time the same model is inconsistent with the data (Supplementary
Section~\ref{S-s8b_evidence:sec:zonesdetail}).
% src: Endo et al. 2021, Results and Discussion (Europe PMC PMC8609560, read 1 October 2026: household 41.3 % of student cases); data/v2_numbers.json (zone.posthoc.schools_Z_ahead_of_H [27, 29]; zone.primary_other.LR_F_vs_Z 42.71, w_hat; zone.TV_boot_ci; zone.check.P_TV_gt_0.15 0.107); notes/VERIFICATION.md (Section 4.9, zone level, items 4, 5, 7); data/v2_numbers.json (zone.seoul: pred_count_median 4, well_mixed_expected 78.7, obs_other 3); data/checks/zones.json: seoul_by_floor (30-120 min gives 7 to 81)

\subsection{What the evidence supports}\label{s8b_evidence:sec:overall}\label{s8_outbreaks:sec:verdict}

\begin{table}[tbp]
  \centering\small
  \caption{The model and its components after both outbreak tests and the analyses of contacts, tracers and
  zones. Every entry is the outcome after the re-check; Supplementary Table~\ref{S-appS_prot2:tab:register}
  registers all tests, with what was pre-specified and what was added post hoc.}
  \label{s8b_evidence:tab:variants}
  % src: the evidence cited is that of Section 7, where each number carries its own src comment; data/v2_numbers.json (out2.pooled.M1.all: delta0 15.446; out2.pooled.M2.all: deltaC -9.150; out2.pooled.plant_share_of_deltaC 0.916)
  \begin{tabular}{@{}>{\raggedright\arraybackslash}p{0.22\textwidth}>{\raggedright\arraybackslash}p{0.48\textwidth}>{\raggedright\arraybackslash}p{0.23\textwidth}@{}}
    \toprule
    Component & Evidence & Position \\
    \midrule
    $\Mone$: people walk on the lattice, transmission on the same site (the model of
      Sections~\ref{s2_model:sec}--\ref{s6_scenes:sec}) &
      expected attack rate below the observed one in three of four seated events of the first test (bus total
      reached with probability 0.0002); in the second test its closed form, which overstates its reach and whose
      pass by the protocol does not bear on the model, beats well mixed ($+15.45$) but fails on three of five new
      events and is worse than a constant ratio; its contacts fail against six of six contact datasets and add
      nothing over a model without space &
      not supported \\\addlinespace[2pt]
    $\Mtwo$: non-local airborne term, coefficient fitted to outbreaks &
      better than well mixed in both tests; the pre-specified pooled comparison favours a constant near/far
      ratio ($-9.15$, nine-tenths from the processing line); tied with a constant two-row ratio on four
      flights; an exponential decay does as well (post hoc); failures on a coach, a call centre, a processing
      line, one flight; no common coefficient; no transfer from trains to aircraft &
      supported only as: a smooth one-parameter gradient beats a well-mixed model \\\addlinespace[2pt]
    $\Mtwo$ with the coefficient measured by tracers &
      one gain over well mixed (classroom), matched by a constant ratio; flight VN54, restaurant and trains
      not reproduced; exploratory extension to four flights ties with a ratio of 3 &
      not validated as a spatial predictor \\\addlinespace[2pt]
    Zone-level reduction, eq.~\eqref{s8b_evidence:eq:zone} &
      weak pass on influenza in 29 schools; a fitted ratio does as well, and a guess of 0.7, 0.1, 0.2 (post hoc)
      does better; building by floor inconclusive &
      weakly supported as structure (class-dominated, three levels); the quantitative law is not \\\addlinespace[2pt]
    Walkers anchored to home sites &
      error of simulated epidemics smaller in a description after the fact; confirmatory test uninformative &
      neither supported nor refuted \\\addlinespace[2pt]
    Heterogeneous mobility $D_x$ &
      no dataset with a measured mobility field; a two-zone walk does not improve on the homogeneous one &
      untested \\
    \bottomrule
  \end{tabular}
\end{table}

The evidence supports the following statement and no stronger one. Tested out of sample against documented
outbreaks, recorded contacts, tracer measurements and a school epidemic, the model does not show general
agreement, and what it gets right is not specific to it (Table~\ref{s8b_evidence:tab:variants},
Supplementary Fig.~\ref{S-s8_outbreaks:fig:summary}). The model with transmission only between people on the same lattice site is
not supported. The extension with a diffusing-aerosol term predicts that risk falls with distance from the source
better than a well-mixed model, with the kernel calibrated on other outbreaks and one infectiousness parameter
fitted per event, but it is not shown to predict the gradient better than a constant ratio of near to far risk:
pooled over the five held-out events of the second test the pre-specified comparison favours the constant ratio,
nine-tenths of it on an event whose near zone was chosen by its source on the same data, and on four held-out
flights the two are statistically indistinguishable. No single mixing coefficient holds across rooms, across
vehicles or from trains to aircraft. Two structural statements are supported, and the model shares them with
simpler ones: risk falls smoothly with distance from the source (pooled near/far ratio with random effects about 4
in the second dataset, 3.3 without the processing line, post hoc; 2.05 over the four held-out events of the first
test with distance strata), and transmission between zones is dominated by the group in which people spend their
time. The contact mechanism, the shape and calibration of the lattice kernel, the reading of fitted coefficients
as a property of the air and their transfer between settings are not supported; attack-rate levels were fitted,
not predicted.
% src: notes/VERIFICATION.md (Section 4.5, first outbreak test, summary); notes/VERIFICATION.md (Section 4.6, second outbreak test, summary); notes/VERIFICATION.md section 4.10 (statement); data/v2_numbers.json (out2.rr_pooled."all 11 events": RR_random 3.978; "without P1" 3.296); data/s8_outbreaks_checks.json (heterogeneity.pooled_rr_random_DL 2.047)

Above all, none of these tests bears on the heterogeneity of mobility, which is what distinguishes the model of
Sections~\ref{s2_model:sec}--\ref{s6_scenes:sec} from simpler ones. In every outbreak that publishes positions,
and hence in every event tested, people are seated or at fixed stations, so the heterogeneity of mobility between
zones never enters, and the model that has a positive result is a standard aerosol model; the contact records
carry no positions; and no metro-carriage or supermarket outbreak with a denominator was found. Nor do the
pre-specified tests bear on the illustrative values $\beta=0.5\perday$ and $\Dz=1\cellday$ of the earlier
sections: in every test the level of transmission was fitted (one infectiousness parameter per outbreak, one rate
for the school epidemic, a matched contact time for the contact records), and mobilities and mixing coefficients
were calibrated on data.
% src: notes/VERIFICATION.md (Section 4.5, first outbreak test, further point 5); notes/VERIFICATION.md section 4.10 (reservations: seated or at fixed stations, mobility untested)
The default values are not meant for single events, and they cannot produce the large ones. Under the
frequency-dependent hazard \eqref{s2_model:eq:pairhazard} an infective at $x$ infects at expected rate $\beta q_x$
(Section~\ref{s2_model:sec:ibm}). Summing, over chains of transmission between distinct people, the expected
number of ordered contact events along each chain (people move independently and remain at stationarity), the
expected number of infections caused by one index case in an event of duration $T$ is at most
$\e^{\beta\qmax T}-1$. With $\beta=0.5\perday$ and the largest $q$ of the scenes ($2.8$) this is $0.10$ for the
1.7-hour Zhejiang bus trip \citep{shen2020}, on which 23 of 67 passengers were infected, and $0.16$ for the 2.5-hour Skagit Valley
choir rehearsal \citep{hamner2020}, at which 52 of 60 attendees (32 of them confirmed) were; by Markov's
inequality the model gives such counts a probability below $0.005$, and the bound reaches them only for
$\beta$ of about $16\perday$ on the bus and $14\perday$ for the choir ($12\perday$ with the 32 confirmed cases).
This comparison was added after the tests, for the two events of the first dataset with the highest attack rates
after a short seated exposure (Supplementary Table~\ref{S-appS_prot2:tab:register}).
% src: data/s8_default_level.json (q_max 2.8; events."Zhejiang bus": expected_infections_bound 0.1021, n_infected 23 of 67, markov_bound 0.0044, beta_needed_per_day 16.34; events."Skagit choir": expected_infections_bound 0.1570, 52 / 32 of 60, markov_bound 0.0030 / 0.0049, beta_needed_per_day 13.61 / 11.99), computed by code/s8_default_level.py from data/bsc_outbreaks/outbreaks.csv
